\documentclass[11pt,letterpaper]{article}
\usepackage[margin=1in]{geometry}
\usepackage{fontspec}
\setmainfont{Latin Modern Roman}
\newfontfamily\chinesefont{Noto Serif CJK SC}
\usepackage{amsmath,amssymb,amsthm,mathtools,booktabs,array,enumitem}
\usepackage[hidelinks,unicode]{hyperref}
\setlength{\emergencystretch}{2em}
\setlist{itemsep=2pt,topsep=5pt}
\newtheorem{theorem}{Theorem}[section]
\newtheorem{proposition}[theorem]{Proposition}
\newtheorem{corollary}[theorem]{Corollary}
\newtheorem{lemma}[theorem]{Lemma}
\theoremstyle{remark}
\newtheorem{remark}[theorem]{Remark}
\DeclareMathOperator{\Aut}{Aut}
\newcommand{\abs}[1]{\lvert #1\rvert}
\newcommand{\ind}{\mathcal I}
\title{Hall bounds under the Mycielski construction\\
and the extremizers of the sixth Mycielski graph}
\author{}
\date{5 October 2026}
\begin{document}
\maketitle
\vspace{-1.5em}
\begin{abstract}
For every nonempty finite simple graph $G$ with $\chi_f(G)<3$, we prove
$\rho(\mu(G))\le3$ and $\rho(\mu^2(G))\le10/3$. The proof uses fractional
coloring, three distinguished vertices, and elementary independence
bounds; no computation enters the upper inequalities. Equality in the
second bound forces an induced subset of order twenty and independence
number six containing all three distinguished vertices. An explicit
fractional coloring of $M_4$ gives $\rho(M_6)\le10/3$; a small exact
certificate for a twenty-vertex witness gives equality. The independently
verified finite census consists of 1,990 maximizing subsets in 199 orbits
under the full ambient group $D_5$. We separate the analytic upper and
equality reductions from the finite witness and census checks. The
strict fractional-coloring hypothesis is a sufficient condition used by
the proof; its necessity is not asserted.
\end{abstract}

\section{Statements and scope}
For a nonempty finite graph $H$, its Hall ratio is
\[
 \rho(H)=\max_{\varnothing\ne S\subseteq V(H)}
 \frac{\abs S}{\alpha(H[S])}.
\]
We abbreviate $\alpha(H[S])$ to $\alpha(S)$ when the ambient graph is clear.
The ordinary Mycielskian $\mu(G)$ contains original vertices $V(G)$, an
independent clone layer $V(G)'$, and an apex $z$. Originals retain their
edges, $u$ is adjacent to $v'$ exactly when $uv\in E(G)$, and the apex
is adjacent to every clone.

\begin{theorem}\label{thm:general}
Let $G$ be a nonempty finite simple graph with $\chi_f(G)<3$. Then
\begin{equation}\label{eq:main}
 \rho(\mu(G))\le3,\qquad \rho(\mu^2(G))\le\frac{10}{3}.
\end{equation}
Any induced subset attaining $10/3$ in $\mu^2(G)$ has order $20$ and
independence number $6$, and contains all three distinguished apex
vertices of the two construction steps. In the four base layers defined
below, it satisfies $p+r=11$ and $q+t=6$. Its final original and clone
cardinalities are one of $(13,6),(14,5),(15,4)$.
\end{theorem}

We use $M_2=K_2$ and $M_{r+1}=\mu(M_r)$, giving orders $2,5,11,23,47$.
An explicit fractional coloring of $M_4$ has weight $29/10<3$. Thus
Theorem~\ref{thm:general} applies to $M_6=\mu^2(M_4)$.

\begin{theorem}[Exact finite specialization]\label{thm:census}
The graph $M_6$ has Hall ratio $10/3$. Exactly $1,990$ labelled subsets
attain this value. They form $199$ orbits under $\Aut(M_6)\cong D_5$,
each of size ten. The orbit counts by final original/clone/apex sizes are
$109$ for $(13,6,1)$, $83$ for $(14,5,1)$, and $7$ for $(15,4,1)$.
\end{theorem}

The upper bound and the necessary equality conditions are analytic.
The matching lower witness and the complete orbit count use exact finite
verification, described in Sections~\ref{sec:M6} and~\ref{sec:census}.
The orbits are those of embedded subsets under the ambient group; they
are not asserted to be distinct abstract graph-isomorphism types.

Cropper, Gy\'arf\'as and Lehel~\cite{CGL} establish the known Hall values
through $M_5$ and unboundedness of the Mycielski sequence. Fractional
coloring under this construction is classical~\cite{LPU}; the Hall-profile
framework also predates this note~\cite{Barnett}. Steiner~\cite{Steiner}
records the separate open question of boundedness of
$\chi_f(M_r)/\rho(M_r)$. The present result does not resolve that
asymptotic question. We make no first-occurrence claim for the general
inequality or the finite census. The hypothesis $\chi_f(G)<3$ is sufficient
for this proof; no necessity or boundary-threshold claim is made.

\section{Fractional coloring and the four-layer model}
A fractional coloring is a collection of nonnegative weights on
independent sets such that every vertex has coverage at least one.
Its total weight bounds the fractional chromatic number. Choose a
fractional coloring of $G$ of total weight $c<3$.

\begin{lemma}\label{lem:projection}
If a graph $X$ admits a homomorphism to $G$, then
\begin{equation}\label{eq:projection}
 \abs{V(X)}\le c\,\alpha(X).
\end{equation}
\end{lemma}
\begin{proof}
Pull back the weighted independent sets along the homomorphism. They
are independent in $X$ and cover every vertex with weight at least one.
Summing the coverage gives
$\abs{V(X)}\le\sum_I w_I\abs I\le c\alpha(X)$.
This also holds when $X$ is empty.
\end{proof}

Write the vertices of $\mu^2(G)$ as four copies of the base and three
special vertices:
\[
 P,\ Q,\ z,\ R,\ T,\ z',\ w.
\]
Here $P$ is the original $G$, $Q$ its first clone layer, and $z$ the first
apex. The second-stage clones of $P,Q,z$ are $R,T,z'$, respectively;
$w$ is the final apex. We use the same letters for the selected vertices
in these layers, with sizes $p,q,r,t$.

Every edge among $P,Q,R,T$ projects to an edge of $G$ by forgetting the
layer. The special adjacencies are: $z$ is adjacent to all of $Q,T$;
$z'$ is adjacent to all of $Q$ and to $w$; and $w$ is adjacent to all of
$R,T,z'$. In particular,
\begin{equation}\label{eq:blocks}
 Q\cup T\text{ is independent},\qquad
 \{z,z'\}\text{ is independent and has no neighbors in }P\cup R.
\end{equation}
These elementary facts supply the refined bound below.

Since $\chi_f(G)<3$, the graph $G$ is triangle-free: a triangle would
require coloring weight at least three. The Mycielski construction
preserves triangle-freeness. A triangle containing the apex is impossible
because its neighborhood is independent; a triangle containing a clone
would project to a triangle in the base.

\section{Proof of the general bounds and equality restrictions}
\subsection{One construction step}
Let $U$ be a nonempty induced subset of $\mu(G)$ with independence number $a$.
Deleting its apex leaves a graph homomorphic to $G$, so
\[
 \abs U\le ca+1<3a+1.
\]
Integrality gives $\abs U\le3a$, proving $\rho(\mu(G))\le3$.

\subsection{Two construction steps}
Let $S$ be a nonempty induced subset of $\mu^2(G)$, set $k=\alpha(S)$, and put
$Z=\{z,z',w\}$. Deleting these three vertices and applying
Lemma~\ref{lem:projection} gives
\begin{equation}\label{eq:large}
 \abs S\le\lfloor ck\rfloor+3\le3k+2.
\end{equation}
For every integer $k\ge6$, $3k+2\le10k/3$, strictly when $k\ge7$.

For $3\le k\le5$, distinguish two cases. If at most two vertices of $Z$
are present, then
\begin{equation}\label{eq:fewapices}
 \abs S\le\lfloor ck\rfloor+2\le3k+1.
\end{equation}
If all three are present, set $X=P\cup R$ and $Y=Q\cup T$.
By~\eqref{eq:blocks}, $\abs Y\le k$ and $\alpha(X)\le k-2$.
Consequently
\begin{equation}\label{eq:allapices}
 \abs S\le\lfloor c(k-2)\rfloor+k+3\le4k-4.
\end{equation}
The resulting bounds for $k=3,4,5$ are $10,13,16$, respectively, all at
most $10k/3$.

The remaining cases use only triangle-freeness. Independence number one
permits at most two vertices. Every triangle-free graph on six vertices
has an independent triple: three neighbors of a vertex form one;
otherwise that vertex has three nonneighbors, among which a nonedge
together with the vertex gives one. Thus independence number two permits
at most five vertices. This proves both inequalities in~\eqref{eq:main}.

\subsection{Equality conditions}
All cases are strict in ratio except possibly $k=3$ or $k=6$.
A triangle-free graph on ten vertices contains an independent four-set:
a vertex of degree at least four supplies one in its neighborhood;
otherwise any vertex has at least six nonneighbors, whose independent
triple can be enlarged by that vertex. Thus $k=3$ cannot attain $10/3$.
Equality requires $k=6$ and $\abs S=20$.

If at most two distinguished vertices were present, the same deletion
argument would give $\abs S\le\lfloor6c\rfloor+2\le19$. Hence all of $z,z',w$ are
selected. Equation~\eqref{eq:blocks} gives
\[
 \abs X\le\lfloor4c\rfloor\le11,\qquad \abs Y\le6.
\]
Since $\abs X+\abs Y=17$, both bounds are tight:
\begin{equation}\label{eq:balance}
 p+r=11,\qquad q+t=6.
\end{equation}
Moreover $\alpha(X)=4$, since $\abs X=11$ and $c<3$, while $Y$ is an
independent six-set.

Let $A$ be the final original part of $S$, an induced subset of $\mu(G)$,
and $B$ its final clone-index set. The selected apex $w$ is nonadjacent to
$A$, so $\alpha(A)\le5$. The one-step result gives $\abs A\le15$.
Since $B'$ is independent, $\abs B\le6$, and
$\abs A+\abs B=19$ implies $\abs A\ge13$. If $\alpha(A)\le4$, then the
same one-step result would give $\abs A\le12$, a contradiction.
Therefore $\alpha(A)=5$ and the three final layer sizes in
Theorem~\ref{thm:general} follow.

\section{The sixth Mycielski graph}\label{sec:M6}
\subsection{An explicit coloring below three}
There is no need to import an exact fractional chromatic value for this
application. Let $C_5$ have its five independent pairs $I$. In
$M_4=\mu(C_5)$, assign the following independent-set weights:
\begin{center}
\begin{tabular}{@{}lcc@{}}
\toprule
Independent sets & Number & Weight per set\\
\midrule
$\{z\}\cup I$ & 5 & $1/5$\\
$I\cup I'$ & 5 & $3/10$\\
The full clone layer & 1 & $2/5$\\
\bottomrule
\end{tabular}
\end{center}
Every cycle vertex belongs to two independent pairs. Thus an original
vertex has coverage $2(1/5+3/10)=1$, a clone has coverage
$2(3/10)+2/5=1$, and the apex has coverage $5(1/5)=1$.
The total weight is $29/10<3$. Theorem~\ref{thm:general} therefore proves
\begin{equation}\label{eq:M6upper}
 \rho(M_6)\le10/3.
\end{equation}
This elementary construction is consistent with the classical fractional
coloring formula of Larsen, Propp and Ullman~\cite{LPU}.

\subsection{A local exact certificate for attainment}
At each construction step, label originals $0,\ldots,n-1$, clones
$n,\ldots,2n-1$, and the apex $2n$. In $M_6$ take
\begin{align*}
 W=\{&0,1,2,3,4,5,6,7,11,12,15,19,21,22,\\
     &26,32,33,41,45,46\}.
\end{align*}
The set $\{11,12,15,19,21,41\}=W\cap(Q\cup T)$ is independent of size six.
An exact finite certificate proves $\alpha(M_6[W])\le6$.

For completeness, the certificate uses only the recurrence
\begin{equation}\label{eq:alpharec}
 a(\varnothing)=0,\qquad
 a(U)=\max\{a(U\setminus\{v\}),\ 1+a(U\setminus N[v])\},
\end{equation}
where $v$ is the least-labelled vertex of nonempty $U$. The two terms
correspond to omitting or including $v$, so induction proves that
$a(U)=\alpha(M_6[U])$. The supplied table contains 190 states, including
the empty state and $W$, and gives $a(W)=6$. Every nonempty row is checked
against its two strictly smaller child states. A separately implemented
include/exclude computation also obtains six. No floating-point
calculation or full $47$-vertex optimization is needed for this witness.

Thus $\abs W/\alpha(W)=20/6=10/3$, proving equality in~\eqref{eq:M6upper}.
The constant in the double-step bound is therefore attained within the
stated class, at $G=M_4$. This attainment does not assert necessity of the
strict hypothesis $\chi_f(G)<3$.

\section{The complete finite census}\label{sec:census}
The analytic equality restrictions reduce the census to final original
sets $A\subseteq V(M_5)$ of orders $13,14,15$ with $\alpha(A)=5$, clone-index
sets $B$ of size $19-\abs A$, and the final apex. Here $N(I)$ denotes the
open neighborhood of $I$ in $M_5$. For such a pair,
\begin{equation}\label{eq:partial}
 \alpha\bigl(M_6[A\cup B'\cup\{w\}]\bigr)
 =\max\left\{1+\alpha(A),
 \max_{\substack{I\subseteq A\\ I\text{ independent}}}
 (\abs I+\abs{B\setminus N(I)})\right\}.
\end{equation}
Indeed, an independent set using $w$ contains no clone; one avoiding $w$
chooses independent originals $I$ and may take all available clones.
Consequently the exact membership criterion is
\begin{equation}\label{eq:membership}
 \abs{B\setminus N(I)}\le6-\abs I
 \quad\text{for every independent }I\subseteq A.
\end{equation}
This criterion does not by itself give a short derivation of the orbit
counts below; those counts come from complete finite enumeration.

\subsection{The ambient group}
The initial $M_3$ is $C_5$, with cyclic order $(0,1,2,4,3)$ in the recursive
labels. At every later stage the apex is the unique maximum-degree
vertex. If the base order is $n$ and its maximum degree is $(n-1)/2$, the
new original degrees are at most $n-1$, clone degrees at most $(n+1)/2$,
and apex degree $n$; this property inducts from $C_5$.

An automorphism fixes the apex and preserves the original and clone
layers. Its original restriction is a base automorphism. The bases are
open-neighborhood twin-free, beginning with $C_5$, so the clone
permutation is uniquely forced by that restriction. Conversely, every
base automorphism extends. Thus $\Aut(M_6)=D_5$, of order ten.
The independent check obtains the ten cycle automorphisms by testing all
120 permutations and then verifies their lifted adjacency preservation.

\subsection{Two independently checked complete enumerations}
The primary program computes $\alpha(A)$ on the $2^{23}$ base subsets
using~\eqref{eq:alpharec}. It retains only the specified original sets
and one representative under $D_5$, leaving 4,370 candidates.
For each $A$, choose an independent five-set $I_0\subseteq A$ minimizing
$\abs{N(I_0)}$. Equation~\eqref{eq:membership} requires at most one clone
index outside $N(I_0)$. Thus enumerate exactly the $b=19-\abs A$ element
sets $B$ consisting of either $b$ elements of $N(I_0)$, or $b-1$ elements
of $N(I_0)$ and one outside. Test all remaining inequalities and
canonicalize the full subset.

This examines 50,504,405 filtered clone candidates. It accepts 200 pairs
at canonical original sets, which reduce to 199 full-subset orbits.
The last deduplication removes the action of an original-set stabilizer
on clone choices.

The independent implementation uses a maximum subset-zeta transform for
the base independence values, and recursive capacity propagation
of~\eqref{eq:membership} for clone choices. Its 184,588 search nodes
produce exactly the same 199 canonical masks. Each pruning rule is
checked for logical sufficiency. Every accepted subset is retested
against the complete independently generated list of 857 maximal
independent sets of $M_6$.

All 1,990 labelled output sets also satisfy the new analytic restrictions:
they contain all three distinguished vertices and have the split
\eqref{eq:balance}. Their stabilizers are trivial. The resulting census is
\begin{center}
\setlength{\tabcolsep}{9pt}
\begin{tabular}{@{}ccc rr@{}}
\toprule
Originals & Clones & Apex & Orbits & Labelled subsets\\
\midrule
13 & 6 & 1 & 109 & 1,090\\
14 & 5 & 1 & 83 & 830\\
15 & 4 & 1 & 7 & 70\\
\midrule
\multicolumn{3}{l}{Total} & 199 & 1,990\\
\bottomrule
\end{tabular}
\end{center}
The analytic reductions prove that no other sizes or layers are omitted;
the exact searches establish the count inside the permitted classes.

\section{Verification boundary and prior work}
The proof of Theorem~\ref{thm:general}, the explicit $29/10$ coloring, and
the necessary equality reductions are analytic. The lower witness has
the small recurrence certificate~\eqref{eq:alpharec}; the census remains
computer-assisted. The earlier 273-node upper-bound certificate and
22-node equality-exclusion certificate are preserved as verification
history, but are no longer premises of the proof.

The source package includes the analytic note, the explicit fractional
coloring, the 190-state witness table and its standard-library checker,
and an independent audit. The earlier census programs, complete orbit
data and independent capacity-propagation check are preserved separately.
No generated format files, caches or page images are needed for the
mathematical reproduction.

The computation uses standard finite recurrences and enumeration.
No improved general maximal-independent-set algorithm is claimed:
Tsukiyama et al.~\cite{Tsukiyama} already give output-sensitive enumeration
in $O(NMK)$ time for a graph with $N$ vertices, $M$ edges and $K$ maximal
independent sets. That row-generation problem is distinct from the exact
Hall optimization and extremizer census.

A bounded current-source check examined the specific
$\chi_f(G)<3$ one-step and two-step Hall statements, including
Barnett's Chapter~4~\cite{Barnett}. It did not locate an equivalent
statement in the inspected material. This is not an exhaustive priority
review. The present wording is ``we prove'', not a first-ever claim.
Neither the general bound nor the finite result resolves the asymptotic
question discussed in~\cite{Steiner}.

\section*{\chinesefont 中文摘要}
\begingroup\chinesefont
\XeTeXlinebreaklocale "zh"
\XeTeXlinebreakskip=0pt plus 1pt
对任意非空有限简单图 $G$，若 $\chi_f(G)<3$，则本文证明
$\rho(\mu(G))\le3$ 及 $\rho(\mu^2(G))\le10/3$。上界证明只使用分数着色、
三个特殊顶点的删除以及初等独立集论证，不依赖有限分支证书。
第二个不等式取等号时，诱导顶点集必有 $20$ 个顶点、独立数 $6$，并包含
全部三个特殊顶点。显式的 $29/10$ 分数着色使该定理适用于 $M_6$；一个
$20$ 顶点见证的 $190$ 状态精确递归证书给出匹配下界。全部 $1\,990$ 个
极值顶点集在环境图的自同构群 $D_5$ 下分成 $199$ 个轨道，轨道计数仍由
两个独立的有限枚举程序验证。严格假设 $\chi_f(G)<3$ 是本证明采用的充分条件，
未断言其必要性或阈值最优性；本文也不作首次发表或渐近问题已解决的声明。
\endgroup

\begin{thebibliography}{9}
\bibitem{CGL}
M.~Cropper, A.~Gy\'arf\'as and J.~Lehel,
\emph{Hall ratio of the Mycielski graphs},
Discrete Mathematics \textbf{306} (2006), 1988--1990.
\href{https://doi.org/10.1016/j.disc.2005.09.020}{doi:10.1016/j.disc.2005.09.020}.
\bibitem{LPU}
M.~Larsen, J.~Propp and D.~Ullman,
\emph{The fractional chromatic number of Mycielski's graphs},
Journal of Graph Theory \textbf{19}(3) (1995), 411--416.
\href{https://doi.org/10.1002/jgt.3190190313}{doi:10.1002/jgt.3190190313}.
\bibitem{Barnett}
J.~B.~Barnett,
\emph{The Fractional Chromatic Number and the Hall Ratio},
Ph.D. dissertation, Auburn University, 2016.
\href{https://etd.auburn.edu/handle/10415/5316}{Auburn dissertation repository},
Definition~2.1.1 and Chapter~4.
\bibitem{Steiner}
R.~Steiner,
\emph{Fractional Chromatic Number Vs. Hall Ratio},
Combinatorica \textbf{45} (2025), article~37.
\href{https://doi.org/10.1007/s00493-025-00164-0}{doi:10.1007/s00493-025-00164-0}.
\bibitem{Tsukiyama}
S.~Tsukiyama, M.~Ide, H.~Ariyoshi and I.~Shirakawa,
\emph{A New Algorithm for Generating All the Maximal Independent Sets},
SIAM Journal on Computing \textbf{6}(3) (1977), 505--517.
\href{https://doi.org/10.1137/0206036}{doi:10.1137/0206036}.
\end{thebibliography}
\end{document}
